
% ============================================================
% SECTION 08
% INVARIANCE AND KNOWLEDGE
% Chapter 02: Invariance
% Book I: Ontology of Viability
% ============================================================

\section{Invariance and Knowledge}
\label{sec:invariance-and-knowledge}

% ------------------------------------------------------------
% 8.1. THE EPISTEMIC PROBLEM OF INVARIANCE
% ------------------------------------------------------------

The analysis of invariance necessarily raises
a question concerning the relationship between
what remains unchanged and the knowledge through
which such invariance is identified.

An invariance claim is not merely an assertion
that something remains unchanged. It also
involves the identification of a determination,
the specification of a domain of transformations,
and the recognition that the determination
persists in a particular respect across that
domain.

These operations involve knowledge.

A knower must distinguish the relevant
determination from other aspects of the
existent, identify or specify the transformations
under consideration, and establish the grounds
for claiming that the determination remains
unchanged.

However, the epistemic conditions under which
invariance is recognized must not be confused
with the ontological status of invariance itself.

The fact that invariance is known through
particular acts of observation, comparison,
or reasoning does not establish that it is
produced by those acts.

The central problem of this section is
therefore to distinguish the ontological
question of invariance from the epistemic
conditions through which invariance becomes
knowable.

% ------------------------------------------------------------
% 8.2. ONTOLOGICAL INVARIANCE AND
% EPISTEMIC RECOGNITION
% ------------------------------------------------------------

The distinction between ontological invariance
and epistemic recognition is fundamental.

Ontological invariance concerns the persistence
of a specified determination in relation to
a domain of transformations.

Epistemic recognition concerns the process
through which a knower identifies, describes,
and justifies the claim that such persistence
occurs.

These two dimensions are related, but they
are not identical.

An existent or one of its determinations
may remain unchanged even when no knower
recognizes that invariance.

Conversely, a knower may identify an apparent
invariance that is not adequately supported
by the relevant conditions or evidence.

The possibility of error therefore belongs
to the epistemic examination of invariance.
It does not imply that the ontological status
of what is being examined is reducible to
the knower's recognition.

The distinction can be summarized as follows:

\begin{enumerate}
    \item Ontological invariance concerns
    what remains unchanged under specified
    transformations.

    \item Epistemic recognition concerns
    how the persistence of that determination
    is identified and justified.
\end{enumerate}

Neither dimension should be substituted
for the other.

% ------------------------------------------------------------
% 8.3. THE CONDITIONS OF RECOGNITION
% ------------------------------------------------------------

The recognition of invariance requires
several analytical operations.

First, the knower must identify the
existent or determinate aspect under
consideration.

Second, the relevant domain of
transformations must be specified.
Without such specification, the
claim of invariance lacks a clear
scope.

Third, the respect in which the
determination is said to remain
unchanged must be identified.

Fourth, the available grounds for
the comparison must be examined.

These operations may involve
observation, measurement, conceptual
analysis, formal reasoning, or
other epistemic procedures.

Their relative importance depends
upon the nature of the existent
and the determination being examined.

The conditions of recognition
are therefore not necessarily
identical across different forms
and domains of invariance.

A quantitative invariance claim
may require different grounds
from a claim concerning a
qualitative, relational, or
structural determination.

This does not imply that one
form of knowledge is universally
superior to another. It indicates
that the adequacy of an epistemic
procedure must be assessed in
relation to the specific claim
being made.

% ------------------------------------------------------------
% 8.4. COMPARISON AND THE IDENTIFICATION
% OF INVARIANCE
% ------------------------------------------------------------

The identification of invariance
often involves comparison between
different conditions or configurations.

A knower examines what is present
under one condition and compares
it with what is present under
another.

The comparison seeks to establish
whether a specified determination
has remained unchanged.

However, comparison does not
automatically establish invariance.

The conditions being compared
may differ in ways that are not
recognized. The descriptions
may omit relevant determinations.
The transformations may not be
adequately specified.

Furthermore, two descriptions
may appear identical while
referring to different ontological
conditions.

The epistemic act of comparison
must therefore be distinguished
from the ontological relation
of invariance that it seeks
to identify.

A comparison is an operation
of the knower. Invariance is
the relation being examined.

The adequacy of the comparison
depends upon whether it captures
the relevant determination and
the transformations within
the specified domain.

% ------------------------------------------------------------
% 8.5. MEASUREMENT, DESCRIPTION,
% AND THE LIMITS OF REPRESENTATION
% ------------------------------------------------------------

In some cases, the identification
of invariance involves measurement
or formal description.

Such procedures can make it
possible to compare determinations
across different conditions
with greater precision.

Nevertheless, the result of
a measurement must not be
identified without qualification
with the complete ontological
status of the existent.

Every measurement operates
through particular procedures,
units, instruments, distinctions,
and assumptions.

A measured quantity may remain
constant within the resolution
and conditions of a procedure,
while other changes remain
unidentified.

Likewise, a formal description
may preserve a particular
relation while omitting other
aspects of the existent.

The limits of measurement and
representation must therefore
be included in the assessment
of an invariance claim.

This does not make measurement
or formalization epistemically
invalid. It establishes that
their results must be interpreted
within the scope of the procedures
through which they are obtained.

The distinction between a
determination and its representation
must remain explicit.

% ------------------------------------------------------------
% 8.6. KNOWLEDGE AND NON-KNOWLEDGE
% ------------------------------------------------------------

The epistemic examination of
invariance takes place within
the situated condition of a
knower.

The knower does not occupy
a position from which every
determination of an existent
is simultaneously available.

Knowledge of one aspect may
coexist with non-knowledge
of other aspects. The identification
of one invariant determination
may occur while other relevant
determinations remain unknown.

These conditions are not
necessarily independent or
sequential.

The knowledge of a particular
domain and the non-knowledge
of other domains may be
intertwined within the same
epistemic situation.

Consequently, the recognition
of invariance must not be
interpreted as the achievement
of complete knowledge of the
existent.

Nor should the limits of
knowledge be treated as evidence
that invariance does not exist.

The absence of an identified
invariant is not equivalent
to the demonstrated absence
of invariance.

Likewise, the recognition of
an invariant in one domain
does not establish that the
existent is invariant in every
other domain.

The epistemic position of the
knower must therefore remain
explicit in every sufficiently
qualified claim about invariance.

% ------------------------------------------------------------
% 8.7. THE DANGER OF EPISTEMIC
% REDUCTION
% ------------------------------------------------------------

A reduction of invariance to
epistemic recognition would
make the existence of invariance
dependent upon whether a knower
has identified it.

Such a position would confuse
the conditions of knowing
with the status of what is
known.

The ontological analysis adopted
in this book does not accept
this reduction.

The existent is not constituted
by the knower's description
of it, and the identity-bearing
character of being is not
exhausted by the knowledge
available to a particular
subject.

Accordingly, invariance must
be examined as a relation
concerning an existent or
one of its determinations,
not merely as a feature of
an epistemic representation.

At the same time, rejecting
epistemic reduction does not
justify ignoring the conditions
and limitations of knowledge.

An ontological claim that
is not epistemically justified
cannot be presented as an
established result merely
because it is asserted to
concern reality.

The analysis must therefore
preserve both dimensions:
the ontological independence
of the existent from a
particular act of knowing,
and the epistemic responsibility
required to make justified
claims about it.

% ------------------------------------------------------------
% 8.8. INVARIANCE, KNOWLEDGE,
% AND THE SPECIFICATION OF DOMAINS
% ------------------------------------------------------------

The specification of a domain
of transformations is both
an analytical requirement
of invariance and an epistemic
requirement for recognizing it.

The domain determines the
range within which the
invariance claim is being
examined.

The knower's ability to identify
that domain, however, may be
limited by the available
knowledge and the conceptual
resources through which
transformations are described.

A domain may be incompletely
specified, or a transformation
may be overlooked.

In such cases, an invariance
claim may be valid within
the transformations actually
examined but unjustifiably
extended beyond them.

The distinction between
a specified domain and an
assumed universal domain
is therefore essential.

The recognition of invariance
within a limited domain does
not establish invariance
under every possible transformation.

This limitation applies even
when the identified invariance
is strongly supported within
the domain that has been
examined.

The epistemic justification
of an invariance claim must
therefore remain proportionate
to the scope of the claim.

% ------------------------------------------------------------
% 8.9. THE STATUS OF INVARIANCE
% AS AN OPEN ONTOLOGICAL QUESTION
% ------------------------------------------------------------

The relational definition of
invariance adopted in this
first edition provides a
framework for distinguishing
what remains unchanged under
specified transformations.

It also makes possible a
disciplined examination of
the relationship between
invariance, change, identity,
continuity, and structure.

However, the adequacy of this
definition as a complete
ontological account has not
been established.

It remains possible that
invariance has a deeper
ontological significance
that cannot be fully expressed
through the relational
definition currently employed.

This possibility must not
be treated as an established
conclusion.

The question is intentionally
left open for later examination,
when the broader ontological
framework has been developed
sufficiently to reconsider
the status of invariance.

The distinction between
ontological invariance and
epistemic recognition remains
necessary regardless of whether
the relational definition is
ultimately retained or revised.

A deeper ontological account
would still need to distinguish
what invariance is from the
conditions under which it
becomes known.

% ------------------------------------------------------------
% 8.10. SYNTHESIS AND TRANSITION
% ------------------------------------------------------------

This section has examined
the relationship between
invariance and knowledge.

The analysis has established
that the ontological status
of invariance must be
distinguished from the
epistemic processes through
which invariance is recognized.

The identification of invariance
involves the specification
of a determination, a domain
of transformations, and the
respect in which something
is said to remain unchanged.

These operations are epistemically
mediated, but they do not
thereby produce the ontological
relation being examined.

The analysis has also emphasized
that knowledge and non-knowledge
are intertwined within the
situated condition of the
knower.

Consequently:

\begin{enumerate}
    \item The absence of knowledge
    of an invariant does not
    establish its ontological
    absence.

    \item The recognition of
    an invariant within one
    domain does not establish
    universal invariance.

    \item The results of
    measurement and description
    must be interpreted within
    the limits of their
    procedures.

    \item The ontological
    independence of invariance
    from a particular act
    of knowing does not
    eliminate the need for
    epistemic justification.

    \item The possibility of
    a deeper ontological
    significance of invariance
    remains open.
\end{enumerate}

The next section examines
the scope and limits of
the concept of invariance.
It will consolidate the
commitments established
throughout this chapter,
distinguish them from
unresolved questions, and
clarify the extent to which
the provisional relational
definition can be used
within the ontology of
viability.

% ============================================================
% END OF SECTION 08
% ============================================================
